Limit-order-book simulators are validated mostly by how realistic their output looks, and used to decide how to execute large orders. We examine whether the first implies the second for three simulator families: ABIDES (agent-based), Microsoft MarS (a generative order-flow model) and a marked multivariate Hawkes order-book simulator. A feasibility audit comes first. We had no licensed order-level data, so realism against real markets was not measured; MarS's released models run on a laptop and respond to injected orders, but cannot be refitted and were unstable, failing in \marsVolFail{} unperturbed runs for its two models. All simulators were then put behind one execution interface with twelve candidate policies and identical cost accounting. In a cross-simulator pilot, policy rankings agreed weakly within a simulator family (mean rank correlation \rankSame{}) and not across families (\rankCross{}). Because disagreement cannot identify who is right, we built four reference markets with known rules, gave each simulator only unperturbed data from them, and measured the cost, in the reference market, of the policy each fitted simulator recommended. Realism did not predict decision quality. The Hawkes fit matched the reference markets' stylised facts best in three of four cases, yet in the two markets where policies differ materially it recommended immediate execution of the larger order, 21 to 22 basis points worse than the best candidate, while a less realistic ABIDES fit was within 2 to 5. Three Hawkes variants with the same realism score recommended different policies. MarS, conditioned on reference order flow, could not be stabilised in any reference market. Averaged over settings, always choosing TWAP (regret 1.2 basis points) beat ABIDES (2.4), a random candidate (4.1) and Hawkes (6.7). The results concern twelve policies in four synthetic markets and are not evidence about real markets.
